% Extra LaTeX preamble for the PDF build (see book.toml in this directory).
% Engine: LuaLaTeX. Chinese text is typeset by luatexja (loaded by pandoc's
% template because `CJKmainfont` is set).

% --- CJK -------------------------------------------------------------------
% luatexja treats Greek/Cyrillic (range 2) and many symbols (range 3: arrows,
% math operators, box drawing, dingbats such as ✨ ➥ ...) as CJK characters by
% default, so they would be taken from the CJK font (full width, often
% missing). Hand them to the Latin fonts and their fallback chains instead;
% this also keeps `├──` trees in code blocks monospaced.
\ltjsetparameter{jacharrange={-2,-3}}
% Without these, luatexja falls back to a Japanese Gothic font for CJK text in
% \textsf / \texttt, which lacks many simplified Chinese characters.
\setsansjfont[JFM=quanjiao]{Noto Sans CJK SC}
\setmonojfont[JFM=quanjiao]{Noto Sans Mono CJK SC}

% --- Code blocks -------------------------------------------------------------
% Wrap long lines instead of letting them run off the page. Pandoc wraps every
% token in a macro (\NormalTok{...}), so breaking inside groups must be allowed.
\usepackage{fvextra}
\DefineVerbatimEnvironment{Highlighting}{Verbatim}{breaklines,breakanywhere,breaknonspaceingroup,commandchars=\\\{\}}
% Light grey background for code blocks (breaks across pages).
\usepackage{framed}
\definecolor{shadecolor}{RGB}{246,248,250}
\renewenvironment{Shaded}{\begin{snugshade}}{\end{snugshade}}

% --- Misc ----------------------------------------------------------------------
% Table of contents: room for three-digit page numbers and for long section
% numbers such as 15.10.32.
\makeatletter
\renewcommand{\@pnumwidth}{2.2em}
\renewcommand{\@tocrmarg}{3.2em}
\renewcommand*\l@section{\@dottedtocline{1}{1.5em}{2.6em}}
\renewcommand*\l@subsection{\@dottedtocline{2}{4.1em}{3.6em}}
\renewcommand*\l@subsubsection{\@dottedtocline{3}{7.7em}{4.4em}}
% Keep chapter titles in running heads as written (no "REDIS"), and let the
% unnumbered introductory chapters update the running head too.
\patchcmd{\chaptermark}{\MakeUppercase}{}{}{}
\apptocmd{\@schapter}{\markboth{#1}{#1}}{}{}
\makeatother
